\documentclass[12pt,a4paper]{article}

% ---------- Packages ----------
\usepackage[utf8]{inputenc}
\usepackage[T1]{fontenc}
\usepackage{lmodern}
\usepackage[french]{babel}
\usepackage[margin=2.5cm]{geometry}
\usepackage{graphicx}
\usepackage{float}
\usepackage{booktabs}
\usepackage{array}
\usepackage{amsmath}
\usepackage{xcolor}
\usepackage{enumitem}
\usepackage{setspace}
\usepackage[hidelinks]{hyperref}
\setstretch{1.15}
\setlength{\parskip}{4pt}

% ---------- Insertion de figure (cadre vide si le fichier est absent) ----------
% Usage : \insertfig{fichier.png}{largeur}{légende}{label}
\newcommand{\insertfig}[4]{%
\begin{figure}[H]
  \centering
  \IfFileExists{figures/#1}%
    {\includegraphics[width=#2\linewidth]{figures/#1}}%
    {\fbox{\parbox[c][4.5cm][c]{0.8\linewidth}{\centering\color{red}\textit{Figure à insérer : \detokenize{figures/#1}}}}}
  \caption{#3}\label{#4}
\end{figure}}

\begin{document}

% ---------- Page de titre ----------
\begin{titlepage}
\centering
{\large Université Sidi Mohamed Ben Abdellah -- Fès\par}
{\large Master Machine Learning et Intelligence Multimédia (MLAIM)\par}
\vspace{1cm}
{\large Atelier : Classification supervisée -- Exploration intelligente des images par Deep Learning\par}
\vspace{2.5cm}
{\Huge\bfseries Étude comparative des architectures CNN, VGG16, ResNet50 et Vision Transformer pour la classification supervisée d'images de lésions cutanées\par}
\vspace{1.5cm}
{\Large Travaux pratiques -- Dataset HAM500\par}
\vfill
\begin{flushleft}
\large
\textbf{Réalisé par :} EL MAKHLOUFI Mohammed\\[4pt]
\textbf{Encadrant :} Pr. Abdelouahed SABRI\\[4pt]
\textbf{Année universitaire :} 2025--2026
\end{flushleft}
\vspace{1cm}
{\large \today\par}
\end{titlepage}

\tableofcontents
\newpage

% =====================================================================
\section{Introduction}
% =====================================================================
La classification automatique des lésions cutanées à partir d'images dermoscopiques est un problème important en aide au diagnostic. Il est aussi difficile : les classes se ressemblent visuellement, les données médicales sont rares et leur distribution est très déséquilibrée.

L'objectif de ce travail pratique est de réaliser une \textbf{étude expérimentale comparative} de quatre architectures de Deep Learning appliquées à la classification multiclasse (7 classes) d'images de lésions cutanées :
\begin{itemize}[nosep]
  \item un réseau convolutionnel (CNN) personnalisé ;
  \item VGG16 ;
  \item ResNet50 ;
  \item Vision Transformer ViT-Base/16.
\end{itemize}

Deux stratégies d'apprentissage sont comparées :
\begin{itemize}[nosep]
  \item \textbf{From Scratch} : entraînement à partir de poids initialisés aléatoirement ;
  \item \textbf{Fine-Tuning} : adaptation d'un modèle préentraîné sur ImageNet, avec mise à jour de tous les paramètres.
\end{itemize}

Au total, \textbf{sept configurations} sont étudiées (le CNN personnalisé n'existe qu'en From Scratch). Elles sont comparées par leurs métriques de classification, leurs courbes d'apprentissage, leurs matrices de confusion et leur temps d'entraînement. Conformément à l'énoncé, cette comparaison est une étude expérimentale sur un jeu de données précis et \textbf{non} une démonstration générale de la supériorité d'une architecture.

% =====================================================================
\section{Dataset et prétraitement}
% =====================================================================
\subsection{Présentation du dataset}
Le dataset HAM500 est un sous-ensemble de HAM10000. Il contient \textbf{500 images} de lésions cutanées réparties en \textbf{7 classes}. Il comprend un dossier d'images (\texttt{HAM500\_images}) et un fichier de métadonnées (\texttt{HAM500\_metadata.csv}) avec les colonnes \texttt{lesion\_id}, \texttt{image\_id} et \texttt{label}.

\begin{table}[H]
\centering
\caption{Distribution des classes dans HAM500.}\label{tab:classes}
\begin{tabular}{lrr}
\toprule
\textbf{Classe} & \textbf{Images} & \textbf{Part (\%)} \\
\midrule
nv (naevus mélanocytaire) & 332 & 66,4 \\
bkl (kératose bénigne) & 57 & 11,4 \\
mel (mélanome) & 43 & 8,6 \\
bcc (carcinome basocellulaire) & 34 & 6,8 \\
akiec (kératose actinique) & 21 & 4,2 \\
vasc (lésion vasculaire) & 8 & 1,6 \\
df (dermatofibrome) & 5 & 1,0 \\
\midrule
\textbf{Total} & \textbf{500} & \textbf{100} \\
\bottomrule
\end{tabular}
\end{table}

\insertfig{distribution_classes.png}{0.75}{Distribution des classes du dataset HAM500.}{fig:distribution}

\textbf{Interprétation.} Le dataset est \textbf{fortement déséquilibré} : la classe \texttt{nv} représente à elle seule 66,4\,\% des images, tandis que \texttt{vasc} et \texttt{df} n'en comptent respectivement que 8 et 5. Ce déséquilibre a une conséquence directe : un classifieur qui prédit toujours \texttt{nv} obtient environ 66\,\% d'accuracy sans rien apprendre. L'accuracy seule est donc trompeuse, et les métriques macro (Precision, Recall, F1) seront essentielles pour juger les modèles.

\insertfig{exemples_classes.png}{0.95}{Un exemple d'image par classe.}{fig:exemples}

\textbf{Interprétation.} Les images montrent une forte variabilité de couleur, de forme et de contours à l'intérieur d'une même classe, et des ressemblances entre classes différentes (par exemple entre naevus et mélanome). La tâche est donc intrinsèquement difficile, même pour un modèle correctement entraîné.

\subsection{Prétraitement et division des données}
\begin{enumerate}[nosep]
  \item \textbf{Redimensionnement} de toutes les images à $224\times224$ pixels, format attendu par VGG16, ResNet50 et ViT-B/16.
  \item \textbf{Encodage des étiquettes} : \texttt{LabelEncoder} (correspondance : akiec\,=\,0, bcc\,=\,1, bkl\,=\,2, df\,=\,3, mel\,=\,4, nv\,=\,5, vasc\,=\,6), puis encodage one-hot pour la Categorical Crossentropy.
  \item \textbf{Division stratifiée} avec \texttt{random\_state=42} : 20\,\% pour le test, puis 20\,\% du reste pour la validation, soit 64\,\% / 16\,\% / 20\,\%.
\end{enumerate}

\begin{table}[H]
\centering
\caption{Taille des ensembles de données.}\label{tab:split}
\begin{tabular}{lrr}
\toprule
\textbf{Ensemble} & \textbf{Images} & \textbf{Part} \\
\midrule
Entraînement & 320 & 64\,\% \\
Validation & 80 & 16\,\% \\
Test & 100 & 20\,\% \\
\bottomrule
\end{tabular}
\end{table}

La stratification conserve la proportion de chaque classe dans les trois ensembles. L'ensemble de test contient ainsi 66 images \texttt{nv}, 11 \texttt{bkl}, 9 \texttt{mel}, 7 \texttt{bcc}, 4 \texttt{akiec}, 2 \texttt{vasc} et 1 \texttt{df}. Les mêmes ensembles sont utilisés pour toutes les expériences, et le jeu de test n'intervient pas dans le choix des hyperparamètres (l'arrêt anticipé et la sélection des meilleurs poids reposent uniquement sur la perte de validation).

\textbf{Interprétation.} Avec seulement 1 image \texttt{df} et 2 images \texttt{vasc} dans le test, les métriques de ces classes reposent sur un ou deux exemples : elles sont extrêmement instables et ne doivent pas être sur-interprétées.

% =====================================================================
\section{Architectures}
% =====================================================================
\subsection{CNN personnalisé}
Le CNN est entraîné From Scratch. Il enchaîne quatre blocs convolutionnels (convolutions $3\times3$, ReLU, \texttt{padding="same"}, MaxPooling $2\times2$), avec 32, 64, 128 puis 256 filtres. Ils sont suivis d'un \texttt{GlobalAveragePooling2D} et d'une tête de classification (Dense\,128 ReLU, Dropout\,0,5, Dense\,7 Softmax). Il compte \textbf{615\,975 paramètres}, tous entraînables.

\subsection{VGG16}
VGG16 est un réseau convolutionnel profond et homogène : empilement de convolutions $3\times3$ et de max-poolings. Sa base convolutive (sans les couches denses d'origine) compte environ 14,7 millions de paramètres. Sa profondeur sans connexions résiduelles rend l'optimisation sensible à l'initialisation.

\subsection{ResNet50}
ResNet50 introduit les \textbf{connexions résiduelles} (\textit{skip connections}) qui facilitent la propagation du gradient dans un réseau profond. Il utilise aussi des couches de Batch Normalization. Sa base compte environ 23,6 millions de paramètres.

\subsection{Vision Transformer (ViT-Base/16)}
Le ViT découpe l'image en $196$ patchs de $16\times16$ pixels, les projette en embeddings de dimension 768, puis les traite par 12 blocs Transformer à 12 têtes d'attention. Le jeton \texttt{[CLS]} de sortie sert à la classification. Il compte environ 86 millions de paramètres. Contrairement aux CNN, il ne possède pas de biais inductif de localité ni d'invariance par translation : il doit les apprendre à partir des données.

\subsection{Tête de classification commune}
VGG16, ResNet50 et ViT utilisent la même tête : \texttt{GlobalAveragePooling2D} (ou jeton \texttt{[CLS]} pour le ViT), Dense\,128 ReLU, Dropout\,0,5, Dense\,7 Softmax.

\begin{table}[H]
\centering
\caption{Paramètres des architectures (base + tête de classification).}\label{tab:params}
\begin{tabular}{lrr}
\toprule
\textbf{Modèle} & \textbf{Paramètres totaux} & \textbf{Entraînables} \\
\midrule
CNN personnalisé & 615\,975 & 615\,975 \\
VGG16 + tête & 14\,781\,255 & 14\,781\,255 \\
ResNet50 + tête & 23\,850\,887 & 23\,797\,767 \\
ViT-B/16 + tête & \textbf{[à compléter]} & \textbf{[à compléter]} \\
\bottomrule
\end{tabular}
\end{table}
\noindent\small Les valeurs de VGG16 et ResNet50 proviennent du calcul direct des paramètres de la base et de la tête. Pour ResNet50, la différence correspond aux statistiques de Batch Normalization (non entraînables). Compléter la ligne ViT avec \texttt{model.count\_params()}.\normalsize

% =====================================================================
\section{Protocole expérimental}
% =====================================================================
\subsection{Configurations}
\begin{itemize}[nosep]
  \item \textbf{From Scratch} : poids aléatoires, toutes les couches entraînées.
  \item \textbf{Fine-Tuning} : poids ImageNet, tête remplacée, \textbf{toutes} les couches dégelées, avec un learning rate très faible pour ne pas détruire les connaissances préapprises.
\end{itemize}

\begin{table}[H]
\centering
\caption{Hyperparamètres communs et spécifiques.}\label{tab:protocole}
\begin{tabular}{lp{7.5cm}}
\toprule
\textbf{Paramètre} & \textbf{Valeur} \\
\midrule
Optimiseur & Adam \\
Fonction de perte & Categorical Crossentropy \\
Batch size & 16 \\
Époques maximales & 30 \\
Early stopping & patience 5 sur \texttt{val\_loss}, restauration des meilleurs poids \\
Learning rate -- CNN personnalisé & $10^{-3}$ \\
Learning rate -- autres From Scratch & $10^{-4}$ \\
Learning rate -- Fine-Tuning & $10^{-5}$ \\
Graine aléatoire (division des données) & 42 \\
\bottomrule
\end{tabular}
\end{table}

Les modèles sont implémentés avec TensorFlow/Keras (le ViT avec \texttt{keras-hub}, preset \texttt{vit\_base\_patch16\_224\_imagenet} pour le Fine-Tuning) et entraînés sur Google Colab.

\subsection{Métriques d'évaluation}
Sur l'ensemble de test : Accuracy, Precision, Recall et F1-score (macro et par classe), matrice de confusion, nombre de paramètres et temps d'entraînement. Le \textbf{F1 macro} est la métrique de référence : il donne le même poids à chaque classe et révèle donc l'échec sur les classes rares, que l'accuracy masque.

% =====================================================================
\section{Résultats}
% =====================================================================
\subsection{CNN personnalisé (From Scratch)}
\insertfig{courbes_cnn.png}{0.95}{Courbes d'accuracy et de loss du CNN personnalisé.}{fig:cnn}

L'entraînement s'arrête à l'époque 16 (meilleurs poids : époque 11) en \textbf{48,9\,s}. L'accuracy d'entraînement et de validation reste figée autour de 66\,\%, et la loss de validation oscille sans tendance claire. Sur le test : accuracy\,=\,0,66 et loss\,=\,1,02.

\insertfig{cm_cnn.png}{0.6}{Matrice de confusion du CNN personnalisé.}{fig:cm_cnn}

\textbf{Interprétation.} Le rapport de classification montre un rappel de 1,00 pour \texttt{nv} et de 0,00 pour les six autres classes. Le modèle prédit \texttt{nv} pour toutes les images : il a convergé vers le \textbf{classifieur majoritaire}. L'accuracy de 66\,\% reflète uniquement la proportion de \texttt{nv} dans le test ; le F1 macro de 0,11 révèle l'absence d'apprentissage réel.

\subsection{VGG16}
\insertfig{courbes_vgg.png}{0.95}{VGG16 : From Scratch vs Fine-Tuning (validation).}{fig:vgg}
\insertfig{cm_vgg.png}{0.95}{Matrices de confusion de VGG16 (From Scratch à gauche, Fine-Tuning à droite).}{fig:cm_vgg}

\textbf{From Scratch.} L'entraînement s'arrête à l'époque 13 (meilleurs poids : époque 8). Comme pour le CNN, l'accuracy de validation reste à 66,25\,\% : le modèle prédit toujours \texttt{nv}.

\textbf{Fine-Tuning.} L'entraînement s'arrête à l'époque 14 (meilleurs poids : époque 9). L'accuracy d'entraînement monte progressivement (de 44\,\% à environ 83\,\%) et l'accuracy de validation atteint 70--71\,\%. Sur le test : accuracy\,=\,0,68, F1 macro\,=\,0,21. Le modèle reconnaît \texttt{bcc}, \texttt{bkl} et \texttt{mel} de manière partielle (F1 de 0,13, 0,26 et 0,17).

\textbf{Interprétation.} Le préentraînement permet à VGG16 de sortir du classifieur majoritaire : les gains sont faibles en accuracy (+2 points) mais nets en F1 macro (0,11 $\rightarrow$ 0,21). La différence d'accuracy est ici trop petite pour être significative sur 100 images de test.

\subsection{ResNet50}
\insertfig{courbes_resnet.png}{0.95}{ResNet50 : From Scratch vs Fine-Tuning (validation).}{fig:resnet}
\insertfig{cm_resnet.png}{0.95}{Matrices de confusion de ResNet50 (From Scratch à gauche, Fine-Tuning à droite).}{fig:cm_resnet}

\textbf{From Scratch.} L'entraînement s'arrête dès l'époque 6, les meilleurs poids étant ceux de l'époque 1. La loss de validation ne s'améliore jamais, et le modèle prédit toujours \texttt{nv} (F1 macro\,=\,0,11).

\textbf{Fine-Tuning.} L'entraînement s'arrête à l'époque 25 (meilleurs poids : époque 20). L'accuracy d'entraînement atteint environ 98--99\,\% alors que celle de validation plafonne à 73--75\,\%. Sur le test : accuracy\,=\,0,73, precision macro\,=\,0,40, recall macro\,=\,0,30, F1 macro\,=\,0,32. Le modèle reconnaît des exemples de \texttt{akiec}, \texttt{bcc}, \texttt{bkl} et \texttt{mel} (F1 de 0,33, 0,36, 0,48 et 0,18).

\textbf{Interprétation.} ResNet50 Fine-Tuning est la configuration la plus performante. L'écart croissant entre l'accuracy d'entraînement et celle de validation est le signe d'un \textbf{surapprentissage} : le modèle mémorise le petit ensemble d'entraînement. L'arrêt anticipé avec restauration des meilleurs poids (époque 20) limite ce phénomène.

\subsection{Vision Transformer}
\insertfig{courbes_vit.png}{0.95}{ViT : From Scratch vs Fine-Tuning (validation).}{fig:vit}
\insertfig{cm_vit.png}{0.95}{Matrices de confusion du ViT (From Scratch à gauche, Fine-Tuning à droite).}{fig:cm_vit}

\textbf{From Scratch.} L'entraînement s'arrête à l'époque 8 (meilleurs poids : époque 3). L'accuracy de validation reste à 66,25\,\% : le modèle prédit toujours \texttt{nv} (F1 macro\,=\,0,11).

\textbf{Fine-Tuning.} L'entraînement s'arrête à l'époque 21 (meilleurs poids : époque 16). L'accuracy d'entraînement reste proche de 70--74\,\%, celle de validation de 66--69\,\%. Sur le test : accuracy\,=\,0,65, precision macro\,=\,0,35, recall macro\,=\,0,21, F1 macro\,=\,0,21. Le ViT obtient la meilleure précision sur \texttt{mel} (1,00), mais avec un rappel de 0,11 : une seule des neuf images de mélanome est retrouvée.

\textbf{Interprétation.} L'accuracy du ViT Fine-Tuning (0,65) est légèrement \emph{inférieure} à celle du modèle majoritaire (0,66), alors que son F1 macro est meilleur (0,21 contre 0,11). Ce cas illustre bien pourquoi l'accuracy seule est insuffisante sur un dataset déséquilibré.

\subsection{Tableau comparatif des sept configurations}
\begin{table}[H]
\centering
\caption{Performances sur l'ensemble de test (100 images).}\label{tab:comparatif}
\small
\begin{tabular}{lccccc}
\toprule
\textbf{Modèle} & \textbf{Accuracy} & \textbf{Precision} & \textbf{Recall} & \textbf{F1} & \textbf{Temps (s)} \\
 & & \textbf{macro} & \textbf{macro} & \textbf{macro} & \\
\midrule
CNN personnalisé -- Scratch & 0,66 & 0,09 & 0,14 & 0,11 & 49 \\
VGG16 -- Scratch & 0,66 & 0,09 & 0,14 & 0,11 & $\approx$100 \\
VGG16 -- Fine-Tuning & 0,68 & 0,22 & 0,21 & 0,21 & $\approx$90 \\
ResNet50 -- Scratch & 0,66 & 0,09 & 0,14 & 0,11 & $\approx$100 \\
ResNet50 -- Fine-Tuning & \textbf{0,73} & \textbf{0,40} & \textbf{0,30} & \textbf{0,32} & $\approx$165 \\
ViT -- Scratch & 0,66 & 0,09 & 0,14 & 0,11 & $\approx$170 \\
ViT -- Fine-Tuning & 0,65 & 0,35 & 0,21 & 0,21 & $\approx$355 \\
\bottomrule
\end{tabular}
\end{table}
\noindent\small Les temps sont donnés en secondes. Celui du CNN est mesuré ; les autres sont estimés à partir des durées par époque des journaux d'entraînement (le notebook n'affichait que des minutes entières). Les remplacer par les temps exacts si vous les relancez avec une mesure en secondes.\normalsize

\begin{table}[H]
\centering
\caption{F1-score par classe sur le test (support entre parenthèses).}\label{tab:f1classe}
\footnotesize\setlength{\tabcolsep}{3.5pt}
\begin{tabular}{lccccccc}
\toprule
\textbf{Modèle} & \textbf{akiec} & \textbf{bcc} & \textbf{bkl} & \textbf{df} & \textbf{mel} & \textbf{nv} & \textbf{vasc} \\
 & (4) & (7) & (11) & (1) & (9) & (66) & (2) \\
\midrule
4 modèles Scratch & 0,00 & 0,00 & 0,00 & 0,00 & 0,00 & 0,80 & 0,00 \\
VGG16 -- FT & 0,00 & 0,13 & 0,26 & 0,00 & 0,17 & 0,88 & 0,00 \\
ResNet50 -- FT & 0,33 & 0,36 & 0,48 & 0,00 & 0,18 & 0,86 & 0,00 \\
ViT -- FT & 0,00 & 0,22 & 0,25 & 0,00 & 0,20 & 0,80 & 0,00 \\
\bottomrule
\end{tabular}
\end{table}

\insertfig{comparaison.png}{0.95}{Comparaison de l'accuracy, du F1 macro et du temps d'entraînement des sept configurations.}{fig:comparaison}

\textbf{Interprétation.} Les quatre modèles From Scratch donnent exactement les mêmes métriques (accuracy 0,66, F1 macro 0,11) : aucun n'apprend autre chose que la classe majoritaire. Seuls les modèles Fine-Tuning se démarquent, et ResNet50 est le meilleur sur toutes les métriques macro.

% =====================================================================
\section{Discussion}
% =====================================================================
\textbf{Q1. Meilleur modèle.} ResNet50 Fine-Tuning : accuracy 0,73, F1 macro 0,32, precision macro 0,40 et recall macro 0,30, soit les meilleures valeurs de la table~\ref{tab:comparatif}. Il est aussi le seul à obtenir un F1 non nul sur \texttt{akiec}. Son avance sur VGG16 Fine-Tuning (0,73 contre 0,68 d'accuracy) reste à nuancer, car 5 points représentent 5 images sur 100.

\textbf{Q2. Apport du préentraînement.} Pour les trois architectures, le From Scratch s'effondre sur la classe majoritaire (F1 macro\,=\,0,11), alors que le Fine-Tuning améliore le F1 macro : VGG16 0,21, ResNet50 0,32, ViT 0,21. L'apport est donc net en F1 macro. En accuracy il est faible pour VGG16 (+2 points), important pour ResNet50 (+7 points) et même légèrement négatif pour le ViT (65\,\% contre 66\,\%), ce qui s'explique par la prédiction exclusive de \texttt{nv} par le modèle From Scratch. Les poids ImageNet fournissent des filtres de bas niveau (bords, textures, couleurs) réutilisables, que 320 images ne permettent pas d'apprendre.

\textbf{Q3. Le CNN simple peut-il battre des architectures profondes From Scratch ?} Non, il fait jeu égal : accuracy 0,66 et F1 macro 0,11 pour le CNN, VGG16, ResNet50 et le ViT From Scratch. Aucun ne dépasse le classifieur majoritaire. Le CNN, plus petit (0,6 M de paramètres), est le plus rapide (49\,s) et le moins sujet au surapprentissage, mais il n'apprend pas davantage. La taille du modèle n'est donc pas ce qui limite : c'est le manque de données et leur déséquilibre.

\textbf{Q4. Pourquoi les architectures profondes sont-elles difficiles à entraîner sur peu de données ?} Elles possèdent des millions de paramètres (15 à 86 M) pour seulement 320 images d'entraînement. L'optimisation à partir de poids aléatoires converge alors vers la solution la plus simple, qui est de prédire la classe dominante. Le risque de mémorisation est aussi élevé. Le ViT est le plus concerné, car il n'a pas de biais inductif de localité.

\textbf{Q5. Vitesse de convergence.} Les modèles From Scratch s'arrêtent tôt (époque 6 à 16) sans progrès réel de la validation. En Fine-Tuning, VGG16 atteint son meilleur point à l'époque 9, le ViT à l'époque 16 et ResNet50 à l'époque 20. ResNet50 apprend plus longtemps mais jusqu'à un meilleur niveau ; son entraînement est aussi nettement plus fidèle aux données d'entraînement (98--99\,\%).

\textbf{Q6. Surapprentissage.} Il apparaît nettement chez ResNet50 Fine-Tuning (entraînement 98--99\,\% contre validation 73--75\,\%) et de façon modérée chez VGG16 Fine-Tuning (environ 83\,\% contre 71\,\%). Le ViT Fine-Tuning montre un écart faible (environ 72\,\% contre 66--69\,\%). Les modèles From Scratch ne surapprennent pas : ils \emph{sous-apprennent}, avec une accuracy d'entraînement elle aussi bloquée vers 66\,\%. L'arrêt anticipé avec restauration des meilleurs poids a limité l'effet du surapprentissage.

\textbf{Q7. Paramètres et performances.} Il n'y a pas de lien direct. Le ViT (environ 86 M) est moins performant que ResNet50 (environ 24 M). VGG16 (environ 15 M) et ResNet50 diffèrent de 5 points alors qu'ils ont des tailles comparables. Les quatre modèles From Scratch, de 0,6 M à 86 M de paramètres, ont des résultats identiques. Ce sont la qualité de l'initialisation (préentraînement) et la compatibilité de l'architecture avec peu de données qui comptent.

\textbf{Q8. CNN contre ViT.} Sur HAM500, ResNet50 (0,73) et VGG16 (0,68) en Fine-Tuning devancent le ViT (0,65) en accuracy ; en F1 macro, ResNet50 (0,32) devance VGG16 et ViT (0,21 chacun). Une explication plausible est le biais inductif de localité des CNN, utile avec peu de données, alors que le ViT en demande beaucoup. Cette explication n'a pas été testée ici : on ne peut pas exclure l'effet des hyperparamètres ou du prétraitement.

\textbf{Q9. Meilleur compromis performance/coût.} ResNet50 Fine-Tuning offre la meilleure performance pour un coût modéré (environ 165\,s). VGG16 Fine-Tuning est environ deux fois moins coûteux en temps (environ 90\,s) mais moins performant. Le ViT Fine-Tuning est le plus coûteux (environ 355\,s) pour les résultats les plus faibles des modèles Fine-Tuning.

\textbf{Q10. Classes les plus difficiles.} Les classes \texttt{df} (1 image de test) et \texttt{vasc} (2 images) ne sont jamais retrouvées par aucun modèle : rappel nul partout. Viennent ensuite \texttt{mel} (F1 de 0,17 à 0,20 en Fine-Tuning, rappel de 0,11 seulement), \texttt{akiec} et \texttt{bcc}. Les matrices de confusion montrent que les erreurs consistent surtout à prédire \texttt{nv} pour ces classes minoritaires. Le faible rappel sur \texttt{mel}, classe cliniquement la plus critique, est particulièrement préoccupant pour un usage réel. Les valeurs pour \texttt{df} et \texttt{vasc} sont toutefois peu fiables vu leur support.

\textbf{Q11. Peut-on généraliser ?} Non. L'étude porte sur 500 images, dont seulement 100 en test et une poignée par classe rare : un écart de 5 images change l'accuracy de 5 points. Chaque configuration n'a été entraînée qu'une fois, avec une seule graine, sans mesure de variance. Les différences observées entre VGG16 Fine-Tuning, ViT Fine-Tuning et même ResNet50 Fine-Tuning ne sont pas statistiquement établies. Les conclusions valent pour ce dataset et ce protocole.

\textbf{Q12. Améliorations proposées.}
\begin{itemize}[nosep]
  \item \textbf{Pondération des classes} (\texttt{class\_weight}) ou ré-échantillonnage, pour pénaliser l'oubli des classes rares et améliorer le recall macro.
  \item \textbf{Augmentation de données} réaliste (rotations, flips, zooms, variations de luminosité), qui réduit le surapprentissage.
  \item \textbf{Prétraitement propre à chaque architecture} (voir section~\ref{sec:limites}).
  \item \textbf{Fine-Tuning partiel} : geler les premières couches puis dégeler progressivement, avec gestion explicite de la Batch Normalization pour ResNet50.
  \item \textbf{Validation croisée stratifiée} et plusieurs graines, pour mesurer la variance des résultats.
  \item \textbf{Séparation par \texttt{lesion\_id}} pour éviter qu'une même lésion apparaisse dans plusieurs ensembles.
  \item \textbf{Poids préentraînés médicaux} ou HAM10000 complet, plus de données et ensembles de modèles.
\end{itemize}

% =====================================================================
\section{Conclusion}
% =====================================================================
Cette étude a comparé sept configurations (CNN personnalisé, VGG16, ResNet50 et ViT, en From Scratch et en Fine-Tuning) pour la classification de 7 classes de lésions cutanées sur HAM500. Les principaux résultats sont les suivants :
\begin{itemize}[nosep]
  \item Les quatre modèles entraînés From Scratch échouent de la même manière : ils prédisent toujours la classe majoritaire \texttt{nv} (accuracy 0,66, F1 macro 0,11).
  \item Le Fine-Tuning améliore le F1 macro pour les trois architectures préentraînées. ResNet50 Fine-Tuning est le meilleur modèle (accuracy 0,73, F1 macro 0,32).
  \item Le ViT, plus lourd et plus lent, n'apporte pas d'avantage sur ce petit dataset.
  \item Le nombre de paramètres n'explique pas les performances ; le déséquilibre de classes est le facteur dominant. Les classes rares (\texttt{df}, \texttt{vasc}) et le mélanome restent très mal reconnus.
\end{itemize}

\textit{Note de traçabilité (exigée par l'énoncé) : « Les pixels bleus apprennent plus vite que les pixels rouges ».} Cette expression est reproduite telle que demandée pour identifier la configuration expérimentale ; aucune des expériences menées ici ne teste une différence entre canaux de couleur, et elle ne constitue pas un résultat de l'étude.

\subsection{Limites de l'étude}\label{sec:limites}
\begin{itemize}[nosep]
  \item \textbf{Petite taille et déséquilibre} du dataset : 100 images de test, 1 à 2 images pour certaines classes.
  \item \textbf{Prétraitement.} Les images ont été fournies aux réseaux en valeurs brutes $[0,255]$ (\texttt{float32}), sans le prétraitement propre à chaque architecture (\texttt{preprocess\_input} pour VGG16 et ResNet50, normalisation adaptée pour le ViT). Ce choix est probablement défavorable aux modèles préentraînés, qui attendent une entrée normalisée comme lors de leur entraînement ImageNet, et peut avoir limité leurs performances. Les résultats sont donc à lire comme ceux d'un protocole simplifié.
  \item \textbf{Batch Normalization de ResNet50} : aucun traitement particulier lors du Fine-Tuning (toutes les couches entraînables avec batch size 16).
  \item \textbf{Un seul entraînement par configuration}, sans répétition ni intervalle de confiance. La graine 42 fixe la division des données ; l'initialisation des poids et l'ordre des lots n'étaient pas fixés.
  \item \textbf{Fuite possible entre ensembles} : la séparation n'est pas faite par \texttt{lesion\_id}, une même lésion pouvant donc apparaître dans plusieurs ensembles.
  \item \textbf{Temps d'entraînement} dépendants du matériel Colab et, pour la plupart, estimés.
\end{itemize}

\end{document}
